\answer{ex:12:1}{$0.60$ (sans corrélation, ce serait $0.18$), limite $\sqrt{c/(rT)}\approx0.58$: cent neurones ne distinguent que \og oui/non\fg{}.}
\answer{ex:12:2}{(a)~$10^8$ impulsions, $\approx10\,\text{mJ}$: $0.3\,\%$ des $3$~J du cerveau entier. (b)~$3$~J, $\approx300$ fois plus.}
\answer{ex:12:3}{(a)~$\theta_A\in[2,3)$, $\theta_B\in[1,2)$; l'autre figure donne au plus $2$ et $1$. (b)~$10\cdot96^2\approx9.2\cdot10^4$.}
\answer{ex:12:4}{(a)~$T_d=0.85\,\tau_a=0.34\,\text{s}$. (b)~$T_d\propto\tau_a$ et ne dépend pas de $I$: $\tau_a\approx3.5\,\text{s}$ (simulation: $3.1\,\text{s}$).}
\answer{ex:12:5}{$p\approx0.17$. La proportion baisse lentement: $0.90$, $0.88$, $0.86$, $0.84$, $0.82$, $0.79$ pour $N=8$, $12$, $24$, $48$, $96$, $192$.}
\answer{ex:12:6}{Les dix essais sont sans erreur pour $\tau_{\text{tr}}=20$, $50$ et $200\ms$ (à $30\ms$, un essai sur dix échoue); à $5$--$10\ms$, environ $0.5$; à $0.5$--$2\,\text{s}$, la qualité tombe à $0.86$. Borne supérieure: la trace doit s'éteindre pendant la pause, $e^{-0.3\,\text{s}/\tau_{\text{tr}}}\ll1$.}
\answer{ex:12:7}{Un seul délai ne suffit pas: la fenêtre de $\pm5\ms$ ne couvre pas $\Delta x/v$ de $5$ à $20\ms$. Deux branches d'inhibition avec $5<d_1\le10\ms$ et $15\le d_2\le d_1+10\ms$.}
\answer{ex:12:8}{\texttt{two\_stage}: $\Delta=2$, égalité dès une inversion, changement de réponse à partir de deux. Compteur de uns: une seule inversion; $0.57$ pour $p=0.1$.}
